\documentclass[11pt,a4paper]{article}

\usepackage[utf8]{inputenc}
\usepackage[T1]{fontenc}
\usepackage{amsmath,amssymb,amsthm,mathtools}
\usepackage{booktabs}
\usepackage{array}
\usepackage{enumitem}
\usepackage{float}
\usepackage{siunitx}
\usepackage{microtype}
\usepackage{lmodern}
\usepackage{listings}
\usepackage{xcolor}
\usepackage[margin=2.4cm]{geometry}
\usepackage[hidelinks]{hyperref}
\usepackage{cleveref}
\usepackage{natbib}

\theoremstyle{plain}
\newtheorem{theorem}{Theorem}[section]
\newtheorem{lemma}[theorem]{Lemma}
\newtheorem{proposition}[theorem]{Proposition}
\newtheorem{corollary}[theorem]{Corollary}
\theoremstyle{definition}
\newtheorem{definition}[theorem]{Definition}
\newtheorem{assumption}[theorem]{Assumption}
\theoremstyle{remark}
\newtheorem{remark}[theorem]{Remark}
\newtheorem{observation}[theorem]{Observation}

\definecolor{codebg}{gray}{0.96}
\lstdefinestyle{sparql}{
  basicstyle=\ttfamily\footnotesize,
  backgroundcolor=\color{codebg},
  frame=single,
  framesep=5pt,
  breaklines=true,
  showstringspaces=false,
  numbers=left,
  numberstyle=\tiny\color{gray},
  numbersep=8pt,
  keywordstyle=\bfseries,
  morekeywords={PREFIX,SELECT,WHERE,VALUES,COUNT,DISTINCT,AS,FILTER,OPTIONAL,UNION},
  commentstyle=\itshape\color{gray},
  morecomment=[l]{\#},
  xleftmargin=1.2em,
}
\lstset{style=sparql}

\newcommand{\Tmap}{\mathcal{T}}
\newcommand{\Adm}{\mathcal{A}}
\newcommand{\Uset}{\mathcal{U}}

\title{\textbf{Database Divergence Is a Measurement, Not a Defect}\\[0.4em]
\large A worked case in federated reaction retrieval,\\
and the conditions under which two answers cannot contradict}

\author{%
Kundai Farai Sachikonye\\[0.3em]
\normalsize Experimental Bioinformatics, TUM School of Life Sciences\\
\normalsize Maximus-von-Imhof-Forum 3, 85354 Freising, Germany\\
\normalsize \texttt{kundai.sachikonye@bitspark.com}}

\date{\today}

\begin{document}
\maketitle

\begin{abstract}
Two SPARQL queries that the specification declares equivalent returned $2$ and
$397$ results respectively against a public biochemical reaction endpoint. The
two differ only in whether a shared property path binds two variables through an
object list (\texttt{?p ?a , ?o}) or through two separate triple patterns. We
document the observation in full, establish by grammar reference that the two
spellings must denote the same abstract query, and confirm on two independent
local engines against a hand-checkable miniature of the same data shape that
both spellings do in fact agree there. The divergence is therefore real,
reproducible against the endpoint, and \emph{not} a specification ambiguity.

The natural conclusion is that a defect exists somewhere in the service stack.
We argue that this conclusion, while defensible, is weaker than it appears, and
that a different reading is available and in our judgement more useful. Under a
propagation model of retrieval in which an answer is the set of items reachable
from a query under a convergent process, we prove that two routes sharing
endpoints and fixed point cannot contradict each other
(\cref{thm:noncontradiction}), and that their numerical difference measures how
much of the underlying network each route resolved
(\cref{thm:extent}). On that reading the pair $(2,397)$ is a measurement of
resolved extent --- specifically, that one route resolved a strict superset of
what the other did --- and reporting it as ``one query is broken'' discards the
quantity it carries.

We are explicit about what this argument does not establish. It requires that
both routes converge to the same fixed point, which for the endpoint in question
we could not verify from outside; without that check the reclassification does
not apply and the ordinary defect reading stands. It also does not establish
that the endpoint conforms to its own documentation, which is a separate
question we do not address. A further confound is documented: the endpoint's
loaded ontology carries $204\,585$ classes against $237\,672$ in the
corresponding download, so the two artefacts are not the same snapshot and no
count comparison across them is clean.

We close with the practical consequence for federated retrieval: a client that
treats divergence between routes as an error condition is discarding
information, and a client that treats it as a measurement gains a coverage
statistic at no additional cost. Neither reading justifies acting on a third
party's service, and we did not.
\end{abstract}

\tableofcontents

\section{Introduction}
\label{sec:intro}

\subsection{The situation}

Ask a public database the same question twice, spelled two ways that the
governing specification says are the same question. You expect the same answer.
When you get $2$ and $397$, three readings are available.

\begin{enumerate}[label=(R\arabic*),leftmargin=3em]
\item \textbf{Specification ambiguity.} The two spellings are not in fact
required to agree, and the expectation was wrong.
\item \textbf{Implementation defect.} The spellings must agree, they do not, and
something in the service stack is broken.
\item \textbf{Extent measurement.} The spellings are two \emph{routes} to the
same query, routes are not part of an answer, and the difference between the
counts measures how much of the data each route reached.
\end{enumerate}

This paper does three things. \Cref{sec:observation,sec:control} establish the
facts well enough to eliminate (R1) --- the specification is unambiguous and two
independent local engines agree with each other and with a hand computation.
\Cref{sec:propagation,sec:reclassification} develop (R3) and prove the two
results that make it a claim rather than a preference.
\Cref{sec:caveats,sec:whatnot} state precisely the conditions under which (R3)
holds, which we could not verify, and are therefore honest that (R2) is not
refuted.

\subsection{Why this is worth a paper rather than a bug report}

Because the choice between (R2) and (R3) is not about this endpoint. Federated
biological retrieval routinely produces multiple routes to nominally the same
question: a live endpoint against a downloaded dump, one ontology version
against another, a property path against an explicit join. Divergences among
those routes are common and are almost universally handled by picking one route
and calling the others wrong. If (R3) is available even sometimes, then that
handling discards a coverage statistic that costs nothing to compute, on every
divergence, in every federated system. The general argument is what we are
after; the observed pair $(2,397)$ is the instance that motivated it.

\subsection{What we did not do}
\label{sec:notdone}

We did not report the observation to the service operator. The finding concerns
a third party's public infrastructure, and the decision whether and how to
report it is the operator's counterpart's to make --- a human one, made by the
person accountable for the relationship, not a side effect of an analysis. We
also ran no load against the endpoint beyond the handful of queries documented
in \cref{tab:runs}, and every mechanistic investigation reported here was
carried out against local engines on synthetic data
(\cref{sec:control}). Nothing in this paper should be read as a recommendation
to probe a public service for the cause.

\section{The observation}
\label{sec:observation}

\subsection{Setting}

The endpoint is a public SPARQL service over a curated resource of biochemical
reactions \citep{Alcantara2012,Bansal2022}, whose participants are identified by
terms from a chemical ontology \citep{Hastings2016}. The question asked is a
standard one for that resource: how many approved reactions have, among their
participants, both a member of a specified amino-acid branch and a member of a
specified oxidant branch, where ``member of a branch'' means transitive
subsumption under one of a small set of ontology roots.

All measurements below were taken on \textbf{10 August 2026} against
\url{https://sparql.rhea-db.org/sparql} with request header
\texttt{Accept: application/sparql-results+json}.

\subsection{The two queries}

The two queries are byte-identical except for one line, reproduced here in full
so the reader need not reconstruct them.

\paragraph{Form A --- object list (comma).}

\begin{lstlisting}[caption={Form A. Returns \textbf{2}.},label=lst:comma]
PREFIX rh:    <http://rdf.rhea-db.org/>
PREFIX rdfs:  <http://www.w3.org/2000/01/rdf-schema#>
PREFIX chebi: <http://purl.obolibrary.org/obo/CHEBI_>
SELECT (COUNT(DISTINCT ?r) AS ?n) WHERE {
  ?r rdfs:subClassOf rh:Reaction ; rh:status rh:Approved .
  ?r rh:side/rh:contains/rh:compound/rh:chebi ?a , ?o .
  VALUES ?aa { chebi:35238 chebi:37022 }
  VALUES ?ox { chebi:35179 chebi:36147 chebi:133294 }
  ?a rdfs:subClassOf* ?aa .
  ?o rdfs:subClassOf* ?ox .
}
\end{lstlisting}

\paragraph{Form B --- two triple patterns (longhand).}

\begin{lstlisting}[caption={Form B. Returns \textbf{397}. Differs from Form A only at lines 6--7.},label=lst:longhand]
PREFIX rh:    <http://rdf.rhea-db.org/>
PREFIX rdfs:  <http://www.w3.org/2000/01/rdf-schema#>
PREFIX chebi: <http://purl.obolibrary.org/obo/CHEBI_>
SELECT (COUNT(DISTINCT ?r) AS ?n) WHERE {
  ?r rdfs:subClassOf rh:Reaction ; rh:status rh:Approved .
  ?r rh:side/rh:contains/rh:compound/rh:chebi ?a .
  ?r rh:side/rh:contains/rh:compound/rh:chebi ?o .
  VALUES ?aa { chebi:35238 chebi:37022 }
  VALUES ?ox { chebi:35179 chebi:36147 chebi:133294 }
  ?a rdfs:subClassOf* ?aa .
  ?o rdfs:subClassOf* ?ox .
}
\end{lstlisting}

\begin{table}[H]
\centering
\small
\caption{Measurements, 10 August 2026.}
\label{tab:runs}
\begin{tabular}{llr}
\toprule
Form & Distinguishing line & \texttt{?n} \\
\midrule
A (object list)      & \texttt{... rh:chebi ?a , ?o .}                          & $2$ \\
B (two patterns)     & \texttt{... rh:chebi ?a .} \; \texttt{... rh:chebi ?o .} & $397$ \\
\bottomrule
\end{tabular}
\end{table}

\subsection{Why the two must denote the same query}
\label{sec:spec}

The comma in a triple block is not an operator. It is punctuation that the
grammar expands before any evaluation semantics is reached.

\begin{observation}[Object lists are syntax]
\label{obs:syntax}
In SPARQL 1.1 \citep{SPARQL11}, a triple pattern with an object list is defined
by the production for \emph{ObjectList} \textup{(}grammar rule 79 in the
\emph{PropertyListPathNotEmpty} family, with \emph{ObjectPath}/\emph{Object}
at rules 83--84 and \emph{GraphNode} forms at 86 and 90\textup{)}. Section
4.2.2 of the specification states directly that an object list abbreviates the
repetition of subject and predicate: the pattern
\texttt{?s ?p ?o1 , ?o2} is defined to expand to
\texttt{?s ?p ?o1 . ?s ?p ?o2 .}
The expansion is performed during translation to the abstract query, before
evaluation. No evaluation rule distinguishes the two.
\end{observation}

\begin{observation}[Property paths are not special-cased]
\label{obs:paths}
The expansion of \cref{obs:syntax} is stated over the predicate slot without
restriction on its form, and the grammar admits a \emph{PathAlternative} in that
slot via \emph{PropertyListPathNotEmpty}. Nothing in the object-list expansion
or in the property-path evaluation semantics (specification §9) makes the
expansion conditional on the predicate being a bare IRI. Hence
\texttt{?r p1/p2/p3/p4 ?a , ?o} expands to the two-pattern form exactly as a
simple predicate would.
\end{observation}

\begin{corollary}
\label{cor:equivalent}
Forms A and B of \cref{lst:comma,lst:longhand} translate to the same abstract
query and must, under a conforming implementation, return the same result.
Reading (R1) of \cref{sec:intro} is therefore eliminated.
\end{corollary}

\subsection{What the correct answer looks like}

It is worth being concrete about the semantics, because the shape of this
query is one where intuition often goes wrong in the direction of expecting the
smaller number.

Both $?a$ and $?o$ range over \emph{the same} set: the ChEBI terms reachable
from reaction $?r$ by the path
\texttt{side/contains/compound/chebi}. They are two independent variables over
one set, not two disjoint slots. A reaction qualifies exactly when that set
contains at least one member under an amino-acid root \emph{and} at least one
member under an oxidant root --- and, notably, a single participant that happens
to lie under both roots would satisfy both conditions simultaneously, since
nothing forces $?a \ne ?o$.

That reading makes $397$ the plausible number and $2$ the surprising one: the
qualifying condition is a conjunction of two existence statements over a set
that is typically of size four to ten, and the branches involved are large. We
note this asymmetry without treating it as proof; the local control of
\cref{sec:control} is what actually settles which value the semantics gives on a
data shape we control.

\section{The local control}
\label{sec:control}

\subsection{Design}

Rather than probe the endpoint further, we reproduced the data \emph{shape} in
miniature: a reaction node linked through the same four-step path to compounds,
with one entity placed under each branch root, small enough that the correct
answer can be computed by hand and checked by inspection.

The miniature contains two reactions:
\begin{itemize}[leftmargin=1.6em]
\item $r_1$ reaches $\{A, O\}$, where $A$ is under an amino-acid root and $O$ is
under an oxidant root;
\item $r_2$ reaches $\{A\}$ only.
\end{itemize}

\paragraph{Hand computation.}
For $r_1$: bind $?a \mapsto A$ (satisfying $A\,\texttt{subClassOf*}\,?aa$) and
$?o \mapsto O$ (satisfying $O\,\texttt{subClassOf*}\,?ox$). Both constraints hold
simultaneously, so $r_1$ qualifies.
For $r_2$: the only value available to either variable is $A$. Binding
$?o \mapsto A$ requires $A$ under an oxidant root, which it is not. No other
binding exists, so $r_2$ does not qualify.
The correct answer is $\{r_1\}$, of size $1$.

\subsection{Result}

Both forms were run unmodified against two independent, unrelated SPARQL
implementations.

\begin{table}[H]
\centering
\small
\caption{Local control. Both engines, both spellings, against the miniature of
\cref{sec:control}. Expected answer by hand computation: $[r_1]$.}
\label{tab:control}
\begin{tabular}{llll}
\toprule
Engine & Form A (comma) & Form B (longhand) & Agreement \\
\midrule
rdflib 7.6.0     & $[r_1]$ & $[r_1]$ & yes \\
pyoxigraph 0.5.9 & $[r_1]$ & $[r_1]$ & yes \\
\midrule
\multicolumn{3}{l}{Hand computation} & $[r_1]$ \\
\bottomrule
\end{tabular}
\end{table}

\begin{observation}[The control is clean]
\label{obs:control}
Two independently developed engines agree with each other, agree across both
spellings, and agree with the hand-computed answer. This confirms
\cref{cor:equivalent} operationally on a data shape matching the endpoint's,
and rules out the possibility that the equivalence fails in practice for
property-path predicates.
\end{observation}

\begin{remark}[What the control does \emph{not} establish]
\label{rem:control-limits}
The miniature has two reactions and a handful of triples. The endpoint has
millions. A control at this scale can establish that the semantics is what we
say it is; it cannot establish anything about behaviour under query planning,
result caps, timeouts, or partial evaluation, all of which are scale phenomena
and none of which the miniature exercises. We built the control to answer one
question --- is the expected answer what we think it is --- and it answers only
that.
\end{remark}

\subsection{A confound: the snapshots differ}
\label{sec:snapshot}

Separately from the query question, the artefacts are not the same data.

\begin{table}[H]
\centering
\small
\caption{Ontology class counts, 10 August 2026.}
\label{tab:snapshot}
\begin{tabular}{lr}
\toprule
Artefact & classes \\
\midrule
Loaded at the SPARQL endpoint            & $204\,585$ \\
Corresponding published download         & $237\,672$ \\
\midrule
Difference                               & $33\,087$ \\
\bottomrule
\end{tabular}
\end{table}

The gap of $33\,087$ classes ($13.9\%$ of the download) means the endpoint's
ontology is a different snapshot, or a filtered subset, or both. Since the
qualifying condition of \cref{lst:comma,lst:longhand} runs through
\texttt{rdfs:subClassOf*} over exactly that ontology, branch membership can
differ between the two artefacts, and \emph{any} count comparison that crosses
them is confounded.

We stress that this confound does not explain the $(2,397)$ pair, because both
forms were run against the same endpoint and therefore the same snapshot. It
matters for a different reason: it means an investigator who tries to adjudicate
the pair by reproducing the query against the download is not running a control,
they are running a third route.

\begin{observation}[Scale context]
For orientation: the resource carries $12\,738$ distinct compounds; the ``has
role'' relation in the lightweight ontology serialisation carries $59\,087$
restrictions; and a representative single metabolite --- a lysine-biosynthesis
intermediate --- participates in exactly $2$ approved reactions. The last figure
is a reminder that per-compound reaction counts in this resource are small,
which is why a jump from $2$ to $397$ at the reaction level is a large relative
change and not a rounding artefact.
\end{observation}

\subsection{Status of the finding}
\label{sec:status}

\begin{center}
\fbox{\parbox{0.92\textwidth}{%
\textbf{Divergence: established.} Two spellings the specification declares
equivalent (\cref{cor:equivalent}), verified equivalent on two independent
engines against a hand-checked miniature (\cref{obs:control}), return $2$ and
$397$ against the same endpoint on the same date.\\[0.5em]
\textbf{Cause: not established.} Query planning, an internal result or
traversal cap, and evaluation-order effects on the shared property path all
remain live candidates. We did not investigate further, for the reason in
\cref{sec:notdone}.}}
\end{center}

\section{A propagation model of retrieval}
\label{sec:propagation}

We now develop the machinery for reading (R3). It is deliberately abstract: the
argument is about routes and answers in general, and the SPARQL case is one
instance.

\subsection{Routes, endpoints, fixed points}

\begin{definition}[Retrieval network]
\label{def:network}
A \emph{retrieval network} is a weighted undirected graph
$(\mathcal{V},\mathcal{E},w)$ whose vertices are addressable items and whose
edge weights $w(e) > 0$ are the cost of traversing a contact between items.
\end{definition}

\begin{assumption}[Contact floor]
\label{ass:floor}
There is $\beta > 0$ with $w(e)\ge\beta$ for every $e\in\mathcal{E}$.
\end{assumption}

\Cref{ass:floor} says no contact is free. For retrieval it is uncontroversial:
following any link costs at least one dereference.

\begin{definition}[Route and admissible set]
\label{def:route}
A \emph{route} $r$ from $v_0$ is a rule for extending contact sequences. Its
\emph{admissible set} at fixed point $x^\ast$ is
\begin{equation}
\Adm_r(v_0, x^\ast) \;=\;
\Bigl\{ v \in \mathcal{V} \;:\;
\exists\, \text{a sequence } v_0 \rightsquigarrow v \text{ permitted by } r
\text{ whose accumulated weight converges to } x^\ast \Bigr\}.
\end{equation}
The \emph{answer} returned by $r$ is $\Adm_r(v_0,x^\ast)$; the query's
\emph{endpoints} are the pair $(v_0, x^\ast)$.
\end{definition}

The two SPARQL forms are two routes in this sense: they start from the same
seed (the reaction class and status constraint), they permit different
extension orders through the shared path, and each returns the set of reactions
it reached.

\subsection{Two structural facts}

\begin{lemma}[Sequences are finite]
\label{lem:finite}
Under \cref{ass:floor}, any contact sequence whose accumulated weight converges
to $x^\ast$ has at most $x^\ast/\beta$ contacts.
\end{lemma}

\begin{proof}
A sum of $n$ terms each at least $\beta$ is at least $n\beta$. Convergence to
$x^\ast$ requires $n\beta \le x^\ast$.
\end{proof}

\begin{theorem}[Inadmissibility has one source]
\label{thm:onesource}
Under \cref{ass:floor}, $v \notin \Adm_r(v_0,x^\ast)$ if and only if no
$r$-permitted sequence from $v_0$ to $v$ converges to $x^\ast$. In particular,
no single mismatched contact along an otherwise convergent sequence can cause
exclusion.
\end{theorem}

\begin{proof}
($\Leftarrow$) By \cref{def:route}. ($\Rightarrow$) Suppose some permitted
sequence converges. By \cref{lem:finite} it is finite, hence exhibits $v$ as a
member. So non-membership forces non-convergence of every permitted sequence.
For the last clause: a mismatched contact can at worst carry larger weight,
which by \cref{lem:finite} shortens the admissible bound rather than removing
the witness; removal requires the whole sequence to fail to converge.
\end{proof}

\begin{theorem}[Route opacity]
\label{thm:opacity}
Let $P,Q$ be distinct convergent sequences from $v_0$ to $v$ reaching the same
$x^\ast$. Then no quantity that is a function of $(v_0,v,x^\ast)$ distinguishes
$P$ from $Q$.
\end{theorem}

\begin{proof}
$P$ and $Q$ agree on all three arguments by hypothesis, so any function of those
arguments takes the same value on both.
\end{proof}

\begin{corollary}[Routes are not part of answers]
\label{cor:notpart}
The answer of \cref{def:route} is defined by \emph{existence} of a convergent
sequence, not by choice of one. Hence the interior of a route is not recoverable
from the answer and was never part of it.
\end{corollary}

\Cref{cor:notpart} is the pivot of the whole argument, and it is worth pausing
on how ordinary it is. A query result is a set of items. It does not contain the
join order, the index chosen, the intermediate bindings enumerated, or the
number of triples touched. Two plans that produce the same set are, to a client
consuming the set, the same answer. Nobody disputes this for query
\emph{plans}; the claim of \cref{sec:reclassification} is that it applies to the
route as a whole and not merely to the plan.

\section{Reclassification}
\label{sec:reclassification}

\subsection{The two results}

\begin{definition}[Resolved extent]
\label{def:extent}
Let $\Uset_r \subseteq \mathcal{E}$ be the set of contacts route $r$ actually
resolved in producing its answer. The \emph{resolved extent} of $r$ is
$|\Uset_r|$; its \emph{unresolved remainder} is
$\mathcal{E}\setminus\Uset_r$.
\end{definition}

\begin{theorem}[Non-contradiction]
\label{thm:noncontradiction}
Let routes $r_1, r_2$ share endpoints $(v_0,x^\ast)$, track the same item set,
and both converge. Then their answers cannot be in contradiction: neither
excludes an item the other admits on grounds of route interior, and any item
excluded by one is excluded by the other on the same grounds
\textup{(}non-convergence\textup{)}.
\end{theorem}

\begin{proof}
By \cref{cor:notpart} the interiors of $r_1$ and $r_2$ are not part of either
answer, so no disagreement about interiors is expressible in the answers. By
\cref{thm:onesource} the only ground for exclusion is non-convergence of every
permitted sequence, which is a property of the pair $(v_0,v)$ and the fixed
point $x^\ast$, shared by hypothesis. Hence the exclusion criteria coincide.
\end{proof}

\begin{theorem}[Divergence measures extent]
\label{thm:extent}
Under the hypotheses of \cref{thm:noncontradiction}, if
$\Adm_{r_1}\ne\Adm_{r_2}$ then $\Uset_{r_1}\ne\Uset_{r_2}$, and the symmetric
difference of the answers is attributable to the symmetric difference of the
resolved extents. Under \cref{ass:floor} each unresolved contact contributes at
least $\beta$ to the corresponding remainder, so
\begin{equation}
\bigl|\,|\Adm_{r_1}| - |\Adm_{r_2}|\,\bigr| \;\le\;
\frac{1}{\beta}\Bigl| \sum_{e\in\Uset_{r_1}} w(e) - \sum_{e\in\Uset_{r_2}} w(e)\Bigr| .
\label{eq:extentbound}
\end{equation}
\end{theorem}

\begin{proof}
By \cref{thm:noncontradiction} the exclusion criteria coincide, so a difference
in answers cannot come from the criteria. The remaining input to
\cref{def:route} that can differ is which permitted sequences were actually
followed, i.e.\ $\Uset_r$. For \cref{eq:extentbound}: each item present in one
answer and absent from the other required at least one contact resolved by one
route and not the other, and each such contact carries weight at least $\beta$
by \cref{ass:floor}; summing gives the bound.
\end{proof}

\begin{corollary}[The pair is a statistic]
\label{cor:statistic}
Under the hypotheses, the observed pair $(2,397)$ states that one route resolved
a strictly larger portion of the network than the other, with the difference in
resolved extent bounded below by $395\beta$. Reporting it as ``one query is
broken'' records the fact of disagreement and discards its magnitude.
\end{corollary}

\subsection{The reading in words}

Stripped of notation: a route is not part of an answer. Two routes that reach
the same place have not disagreed about anything the answer contains; they have
covered different amounts of ground getting there. The number $397 - 2 = 395$ is
then not an error magnitude --- it does not say ``$395$ results were lost'' in
the sense of a fault --- it says the two routes differ in coverage by at least
that much. That is a quantity about the \emph{retrieval}, and it is available
for free whenever a system happens to run two routes.

\subsection{Why the alternative reading is not thereby refuted}

\Cref{cor:statistic} is conditional, and the condition is not decorative.
\Cref{thm:noncontradiction} requires both routes to converge to the same
$x^\ast$. If they do not --- if the two spellings are, at the level of the
engine's execution, asking different questions of different portions of the
store --- then the theorem's hypothesis fails, the routes are not two routes to
one query, and disagreement between them means exactly what it appears to mean.

We could not verify convergence to a common fixed point from outside the
service. Verifying it would require observing the execution, which is precisely
the investigation we declined to run against a third party's infrastructure
(\cref{sec:notdone}). So the honest statement is: (R3) is available if the
convergence condition holds, we did not establish that it holds, and therefore
(R2) --- an implementation defect --- remains a live and reasonable reading of
the same data.

\section{Caveats, stated as such}
\label{sec:caveats}

\paragraph{Caveat 1: convergence to a common fixed point is unverified.}
This is the load-bearing one. Every result in
\cref{sec:reclassification} assumes it. Without it,
\cref{thm:noncontradiction,thm:extent,cor:statistic} do not apply to the
observed pair. A reader who wants the reclassification must either verify the
condition or treat the conclusion as conditional; we treat it as conditional.

\paragraph{Caveat 2: nothing here validates the endpoint against its documentation.}
Whether the service behaves as its own documentation says is a separate question
from whether two routes through it can contradict. We have not examined the
documentation and make no claim about conformance. In particular, if the service
documents a result cap or a traversal bound, that documentation may fully
explain the observation, and our argument neither anticipates nor contradicts
it.

\paragraph{Caveat 3: the control is small.}
Per \cref{rem:control-limits}, the miniature establishes semantics and nothing
about scale behaviour. Since scale behaviour is the leading candidate cause, the
control is silent on cause by construction.

\paragraph{Caveat 4: the snapshots differ.}
Per \cref{sec:snapshot}, endpoint and download carry different ontology class
counts. Any attempt to adjudicate by cross-artefact comparison introduces a
third route rather than a control.

\paragraph{Caveat 5: single date, single endpoint.}
All measurements are from one date against one service. We did not test
persistence across days, alternative mirrors, or the same query shape against
other resources. A divergence that is not reproducible over time is a different
phenomenon from one that is, and we do not know which this is.

\section{What follows for federated retrieval}
\label{sec:practice}

Independently of which reading is right for this case, the framework has a
consequence that costs nothing to adopt.

\begin{enumerate}[label=(P\arabic*),leftmargin=3em]
\item \textbf{Record both answers.} A federated client that runs two routes and
keeps only one has thrown away the pair. Keeping both costs one integer.
\item \textbf{Report the difference as coverage, not as error.} By
\cref{thm:extent} the difference bounds a resolved-extent gap. Even if the
convergence hypothesis is unverified, ``routes A and B differ by $395$'' is a
strictly more informative log line than ``route A failed''.
\item \textbf{Check the convergence hypothesis before reclassifying.} The
reclassification is conditional. A system that applies it unconditionally will
silently absorb genuine faults as coverage differences, which is the failure
mode this framework introduces and the reason (P4) exists.
\item \textbf{Keep a route that can fail.} If every route in a federation is
tolerant enough to converge on anything, divergence stops being informative.
At least one route should be strict enough that its disagreement is meaningful.
\item \textbf{Never compare counts across snapshots.} Per \cref{sec:snapshot}.
This one is not conditional on anything.
\end{enumerate}

\section{What this paper does not claim}
\label{sec:whatnot}

\begin{itemize}[leftmargin=1.6em]
\item That the endpoint is correct. We do not know.
\item That the endpoint is defective. We do not know that either; we established
divergence, not cause (\cref{sec:status}).
\item That $397$ is the right answer and $2$ the wrong one. The semantic
argument of \cref{sec:spec} makes $397$ the more plausible value on a full
store, and the local control shows the semantics behaves as expected in
miniature, but neither computes the correct value on the endpoint's actual data.
\item That the reclassification of \cref{sec:reclassification} applies to this
case. It applies if the convergence hypothesis holds, which is unverified
(Caveat 1).
\item That divergences in general should not be investigated as bugs. Sometimes
they are bugs. The claim is that ``bug'' and ``coverage measurement'' are
distinguishable by a checkable condition, and that the check is worth running
before choosing.
\end{itemize}

\section{Conclusion}

Two spellings that SPARQL 1.1 §4.2.2 defines as the same abstract query, and
that two independent engines confirm behave identically on a hand-checkable
replica of the data shape, returned $2$ and $397$ against a public endpoint on
10 August 2026. The divergence is established; its cause is not, and we declined
to pursue the cause against a third party's service.

The conventional reading is that something is broken. We have set out an
alternative under which the observation is a measurement: routes are not part of
answers (\cref{cor:notpart}), two routes sharing endpoints and fixed point
cannot contradict each other (\cref{thm:noncontradiction}), and their difference
therefore quantifies resolved extent rather than error
(\cref{thm:extent,cor:statistic}). The alternative is conditional on a
convergence check we could not perform from outside, and we have said so
plainly rather than presenting a conditional result as a settled one.

What survives either reading is the practice in \cref{sec:practice}: a federated
client that logs the divergence as a magnitude rather than a failure is strictly
better informed, and it costs one integer. The framework's own risk --- that
applying the reclassification without the check would absorb real faults as
coverage differences --- is stated as (P3) precisely because it is the thing that
would make this argument harmful if adopted carelessly.

\bibliographystyle{plainnat}
\bibliography{references}

\end{document}
